% Unique practice relocated from topic 2.1; no new board/year attribution.
\begingroup
\chTwoReadableTopicStyle
\subsection{যোগাযোগ ও ব্যান্ডউইথ: স্থানান্তরিত অনুশীলন}
এগুলো নিজস্ব ধারণাভিত্তিক অনুশীলন। মূল পাঠের বড় প্রশ্নতালিকা থেকে আলাদা করা হয়েছে; একই প্রশ্ন আবার মূল পাঠে নেই।
\textbf{MCQ}
\begin{multicols}{2}\begin{enumerate}
\item\chTwoMCQBlock{কমিউনিকেশন সিস্টেমে কোন অংশ ডেটা পাঠায়?
\par ক. প্রাপক \quad খ. প্রেরক \par গ. মাধ্যম \quad ঘ. নয়েজ }
\item\chTwoMCQBlock{বাইনারি ডিজিটাল সংকেতের দুটি মান কী?
\par ক. A ও B \quad খ. 1 ও 2 \par গ. 0 ও 1 \quad ঘ. X ও Y }
\item\chTwoMCQBlock{ডিজিটাল সংযোগের বহনক্ষমতা হিসেবে ব্যান্ডউইথের একক কোনটি?
\par ক. মিটার \quad খ. bps \par গ. ভোল্ট \quad ঘ. ওয়াট }
\item\chTwoMCQBlock{8000 bit ডেটা 4 সেকেন্ডে পাঠালে গড় ডেটা হার কত?
\par ক. 500 bps \quad খ. 1000 bps \par গ. 2000 bps \quad ঘ. 4000 bps }
\item\chTwoMCQBlock{অ্যানালগ সংকেতের মানের প্রকৃতি কেমন?
\par ক. বিচ্ছিন্ন \quad খ. অবিচ্ছিন্ন \par গ. শুধু বাইনারি \quad ঘ. প্যাকেটভিত্তিক }
\item\chTwoMCQBlock{প্রচলিত পরিষেবা-ছকে স্ট্রিমিং ভিডিওর উদাহরণ ডেটা হার কত?
\par ক. 0.5 Mbps \quad খ. 0.7 Mbps \par গ. 1.5 Mbps \quad ঘ. 4 Mbps }
\end{enumerate}\end{multicols}
\textbf{উত্তরমালা:} ১. খ, ২. গ, ৩. খ, ৪. গ, ৫. খ, ৬. খ।

\textbf{CQ: অনলাইন ক্লাসের মান (নিজস্ব অনুশীলন)}\par
এক শিক্ষার্থী গ্রামের বাড়ি থেকে 4G সংযোগে অনলাইন ক্লাস করছে। ভিডিও মাঝে মাঝে আটকে যায়, শব্দ দেরিতে আসে এবং ফাইল পাঠাতে সময় লাগে।
\begin{enumerate}
\item ক. Bandwidth কী?
\item খ. কম্পিউটার যোগাযোগে বাইনারি ডিজিটাল সংকেত উপযোগী কেন?
\item গ. উদ্দীপকের সমস্যার সঙ্গে প্রকৃত ডেটা হার বা Throughput এবং বিলম্ব বা Latency-এর সম্পর্ক ব্যাখ্যা করো।
\item ঘ. শুধু 4G সংযোগ থাকলেই ভালো অনলাইন ক্লাস নিশ্চিত হয় কি? যুক্তিসহ বিশ্লেষণ করো।
\end{enumerate}
\endgroup
